\section{Теорема о том, что сумма собственных подпространств является прямой. Критерий существования базиса из собственных векторов}

\begin{theorem}[О прямой сумме собственных подпространств]
Пусть $\lambda_1, \lambda_2, \dots, \lambda_k$ --- попарно различные собственные числа линейного оператора $\varphi$. Тогда сумма соответствующих собственных подпространств является прямой:
\[
L_{\lambda_1} + L_{\lambda_2} + \dots + L_{\lambda_k} = L_{\lambda_1} \oplus L_{\lambda_2} \oplus \dots \oplus L_{\lambda_k}.
\]
\end{theorem}

\begin{proof}
Для доказательства того, что сумма подпространств прямая, достаточно показать, что если $x_1 + x_2 + \dots + x_k = 0$, где $x_i \in L_{\lambda_i}$, то $x_i = 0$ для всех $i = 1, \dots, k$.

Предположим, что не все $x_i$ равны нулю. Тогда мы получили бы нетривиальную линейную комбинацию векторов $x_1, \dots, x_k$, равную нулю. Однако, по основной теореме о собственных векторах (билет 30), система собственных векторов, соответствующих попарно различным собственным числам, всегда линейно независима. 
Это означает, что равенство $x_1 + \dots + x_k = 0$ возможно только в тривиальном случае, когда $x_1 = x_2 = \dots = x_k = 0$. 
Следовательно, пересечение любого из подпространств $L_{\lambda_i}$ с суммой остальных равно $\{0\}$, что и означает, что сумма является прямой.
\hfill$\blacksquare$
\end{proof}

\begin{theorem}[Критерий существования базиса из собственных векторов]
Линейный оператор $\varphi$ в $n$-мерном пространстве $V$ имеет базис, состоящий из собственных векторов, тогда и только тогда, когда выполнены два условия:
\begin{enumerate}
    \item Характеристический многочлен $P_\varphi(\lambda)$ распадается на линейные множители над полем $F$ (то есть имеет ровно $n$ корней с учетом их алгебраических кратностей).
    \item Для каждого собственного числа $\lambda_i$ его геометрическая кратность равна алгебраической: $m_g(\lambda_i) = m_a(\lambda_i)$.
\end{enumerate}
\end{theorem}

\begin{proof}
\begin{itemize}
    \item \textbf{Необходимость ($\Rightarrow$):} Пусть существует базис из собственных векторов. Тогда матрица оператора $A_\varphi$ в этом базисе диагональна: $A_\varphi = \operatorname{diag}(\lambda_1, \dots, \lambda_n)$. 
    Характеристический многочлен имеет вид $P_\varphi(\lambda) = (\lambda_1 - \lambda)(\lambda_2 - \lambda)\dots(\lambda_n - \lambda)$, что доказывает условие 1. 
    Если некоторое число $\lambda^*$ встречается на диагонали $m_a$ раз, то система $(A_\varphi - \lambda^* E)X = 0$ будет иметь ровно $n - m_a$ нулевых строк после приведения к ступенчатому виду, значит, размерность пространства решений (геометрическая кратность) равна $m_a$. Условие 2 выполнено.
    
    \item \textbf{Достаточность ($\Leftarrow$):} Пусть выполнены условия 1 и 2. Пусть различные собственные числа --- это $\lambda_1, \dots, \lambda_s$. 
    По условию 1, сумма их алгебраических кратностей равна размерности пространства: $\sum_{i=1}^s m_a(\lambda_i) = n$.
    По условию 2, $m_g(\lambda_i) = m_a(\lambda_i)$ для всех $i$. Следовательно, $\sum_{i=1}^s m_g(\lambda_i) = n$.
    
    Рассмотрим сумму собственных подпространств $W = L_{\lambda_1} + \dots + L_{\lambda_s}$. По предыдущей теореме, эта сумма прямая. Размерность прямой суммы равна сумме размерностей слагаемых:
    \[
    \dim W = \sum_{i=1}^s \dim L_{\lambda_i} = \sum_{i=1}^s m_g(\lambda_i) = n.
    \]
    Поскольку $W$ --- подпространство $n$-мерного пространства $V$ и $\dim W = n$, то $W = V$. 
    Если в каждом $L_{\lambda_i}$ выбрать базис (он будет состоять из $m_g(\lambda_i)$ векторов), то объединение этих базисов даст систему из $n$ векторов. Поскольку сумма прямая, эта объединенная система линейно независима (билет 18). 
    Система из $n$ линейно независимых векторов в $n$-мерном пространстве образует базис. Все эти векторы являются собственными.
\end{itemize}
\hfill$\blacksquare$
\end{proof}
